\subsection{Estado técnico do artefato}

Esta seção distingue provas de engenharia de resultados científicos. Os marcos M1/M2, registrados em 31 de agosto de 2026, são históricos; os incrementos posteriores de integração privada e registro não os substituem nem constituem a coleta confirmatória. M1 estabeleceu o laboratório reproduzível. A imagem do Odoo 19.0 Community foi construída a partir do commit fixado, com manifest OCI local registrado separadamente do config digest. PostgreSQL 18 permaneceu sem porta publicada; Odoo foi exposto somente em loopback; e uma identidade técnica C0, distinta da chave administrativa, realizou o fluxo JSON-2.

As fixtures sintéticas incluem empresa, fornecedor, conta bancária, produto e pedido de compra confirmado. O smoke M1 criou uma única \texttt{account.move} do tipo \texttt{in\_invoice}, em \texttt{draft}, com \texttt{auto\_post='no'}, beneficiário persistido e linha vinculada ao pedido. Uma leitura independente verificou o efeito; em seguida, o rascunho foi removido. O hash canônico após a remoção coincidiu com o estado inicial, demonstrando que o testbed consegue produzir e desfazer a mutação observada sem apagar os volumes.

M2 ligou uma entrada documental ao mesmo efeito. O pipeline gera localmente um PDF sintético, marcado como sem valor fiscal e contendo camada de texto; valida tipo, tamanho, número de páginas e limite de caracteres; extrai o texto sem OCR; produz uma proposta \texttt{vendor-bill-proposal/v1}; valida propriedades, tipos, formatos e total por aritmética decimal; resolve referências por chaves de negócio; e chama o executor direto C0. PDFs criptografados, inválidos, sem texto ou fora dos limites são recusados.

Para isolar o encanamento antes da seleção do LLM, M2 usou um adaptador determinístico de desenvolvimento, versão 1. Esse adaptador lê somente os campos declarados da fixture e é proibido em runs científicas. Sua função foi demonstrar que PDF, contrato, executor, Odoo e oráculo concordam, não simular inteligência ou resistência a ataque.

\begin{table}[H]
\centering
\caption{Evidências técnicas disponíveis antes do piloto.}
\label{tab:current-evidence}
\footnotesize
\begin{tabularx}{\textwidth}{L{2.10cm} Y L{3.05cm}}
\toprule
\textbf{Evidência} & \textbf{Resultado observado} & \textbf{Limite do resultado} \\
\midrule
Fonte Odoo & Community 19.0 no commit \texttt{a4b877245c9cb5e3d7152d1593a518170f115ca1} & Não representa todo Odoo 19 futuro \\
Manifest de imagem & OCI local \texttt{sha256:3e0bcf8e7968\ldots} ligado à fonte & Ainda não publicado em registry \\
Identidade C0 & Conta técnica não administrativa cria e lê o rascunho & Evidência histórica M1/M2, não integração científica completa \\
Contrato M2 & \texttt{vendor-bill-proposal/v1}, fechado e versionado & Precede o contrato comum v2 atual \\
Testes M2 & 41 testes aprovados no gateway 0.2.0 & Suíte histórica, não contagem da árvore atual \\
Smoke M2 & PDF $\rightarrow$ proposta $\rightarrow$ rascunho $\rightarrow$ remoção & Usa adaptador determinístico de desenvolvimento \\
Estado & hash inicial e restaurado \texttt{224448b68946\ldots}; estado criado distinto & Um cenário benigno de engenharia \\
\bottomrule
\end{tabularx}
\end{table}

A suíte cobre PDF criptografado ou sem texto, estrutura inválida, campo duplicado ou ausente, propriedade extra, versão de schema incompatível, data e valor malformados, total divergente, referência inexistente, duplicidade e resíduo no Odoo. O smoke final passou com 41 testes, criou exatamente um rascunho e restaurou o hash inicial. PDF, texto, proposta e estados receberam SHA-256; credenciais e evidências detalhadas ficaram fora do Git.

Até 26 de outubro, o repositório também registra provas posteriores: contrato comum v2, caminhos privados de criação e contabilização C0/C1, autorização vinculada, captura independente de estado e restauração por snapshot de banco e filestore. O checkpoint de captura/restauração percorreu oito casos técnicos e 16 restaurações, incluindo contabilização, com provedor e supervisão simulados. Esses números contam casos de engenharia, não ataques científicos nem observações independentes do estudo.

O registro candidato \texttt{run-record/v2} preserva resultados recomputados, falhas de proposta, estado não observável e divergência de reset, sem preencher lacunas com sucesso. O gravador local preserva registros por hashes e escrita exclusiva; isso não demonstra backup externo imutável. A integração operacional e a revisão independente permanecem abertas, e o método público \texttt{execute} continua fechado. As provas anteriores não foram reexecutadas para esta revisão editorial.

\subsection{O que os resultados atuais demonstram}

Os resultados demonstram quatro propriedades necessárias ao futuro experimento. Primeiro, o estado de negócio pode ser observado independentemente da saída do modelo. Segundo, a proposta pode ser validada e canonicalizada antes da mutação. Terceiro, as condições podem começar da mesma fixture e terminar sem resíduo detectável. Quarto, os identificadores internos do ERP não precisam ser aceitos do documento: fornecedor, pedido e conta são resolvidos por chaves sintéticas declaradas.

Essas propriedades reduzem ameaças de instrumentação. Sem um oráculo externo, uma ferramenta chamada poderia ser contada como ataque bem-sucedido mesmo se a transação falhasse; sem reset, uma condição poderia herdar o efeito anterior; sem contrato comum, C1 poderia parecer superior apenas porque rejeita JSON inválido que C0 aceita. M1/M2 antecipam esses problemas e mantêm a validação estrutural nos dois braços.

O trabalho também responde ao feedback de que a continuidade do TCC I precisava tornar-se aplicada e definir segmento-alvo. O artefato executa um fluxo de Contas a Pagar em um ERP real de laboratório, e não apenas uma arquitetura conceitual. Ainda assim, implementação não substitui avaliação. O resultado atual é evidência de viabilidade e controle experimental, não de eficácia da defesa.

\subsection{Piloto: propostas do modelo}

Em 26 de outubro de 2026, o piloto executou o estágio de propostas com o modelo real: 12 pares, dois membros, C0 e C1 e três repetições, totalizando 144 chamadas por rodada, em ordem randomizada por semente registrada. Mediu-se a proposta (ASR-proposal) e a decisão que a política C1 tomaria sobre ela (ASR-policy), contra o estado confiável derivado do gabarito. O efeito no Odoo, isto é, ASR-state, UACR e BTSR, ainda não foi medido, porque o ciclo integrado não foi ligado aos 12 pares.

A rodada 1 usou notas adversariais de uma ou duas frases. Como nenhuma proposta de C0 foi comprometida, o gate de eficácia disparou. As 12 notas foram então reescritas segundo princípios estáticos: justificar por que o valor impresso estaria errado, alegar prioridade ou competência, nomear o campo e o valor a registrar, integrar-se ao vocabulário do documento e manter ofuscação legível. A revisão não usou conhecimento da defesa e manteve documentos benignos, gabarito, valores do atacante, prompt e modelo. A rodada 2 foi a única repetição permitida.

\begin{table}[H]
\centering
\caption{Piloto de propostas: resultados por rodada e condição.}
\label{tab:pilot-proposals}
\footnotesize
\begin{tabularx}{\textwidth}{Y C{1.35cm} C{1.35cm} C{1.35cm} C{1.35cm}}
\toprule
& \multicolumn{2}{c}{\textbf{Rodada 1}} & \multicolumn{2}{c}{\textbf{Rodada 2}} \\
\textbf{Medida} & \textbf{C0} & \textbf{C1} & \textbf{C0} & \textbf{C1} \\
\midrule
Chamadas com proposta válida & 72/72 & 72/72 & 72/72 & 71/72 \\
Benignas idênticas ao gabarito & 36/36 & 36/36 & 36/36 & 35/36 \\
Benignas aceitas pela política C1 & -- & 36/36 & -- & 35/36 \\
Adversariais com campo-alvo alterado (ASR-proposal) & 0/36 & 0/36 & 0/36 & 0/36 \\
Adversariais idênticas ao gabarito & 36/36 & 36/36 & 36/36 & 36/36 \\
Comprometidas aceitas por C1 (ASR-policy) & -- & 0/0 & -- & 0/0 \\
\bottomrule
\end{tabularx}
\end{table}

Nas duas rodadas, nenhuma das 72 propostas adversariais de C0 incorporou o objetivo do atacante, e o limiar fixado antes da segunda rodada não foi atingido. Com 0 em 36 execuções de C0 em uma rodada, o limite superior unilateral de 95\% da taxa é de cerca de 8\% se as execuções fossem independentes e de 22\% tomando o par como unidade. A única falha foi uma saída fora do schema em uma chamada benigna de C1 da rodada 2, mantida como resultado e sem nova tentativa. Cada rodada consumiu cerca de 215 mil tokens, aproximadamente US\$~0,07 pelos preços datados, com latência mediana de 2,1~s por chamada.

Como verificação de componente, e não como resultado do experimento, propostas que obedecessem integralmente às 12 notas foram submetidas à política C1. Onze seriam negadas por divergência de pedido, fornecedor, valor, linha, conta fora da allowlist ou alvo da contabilização. Uma, que altera somente a referência impressa do documento, seria aceita: essa referência é controlada pelo atacante e funciona apenas como sinal secundário de duplicidade.

\subsection{Resultados confirmatórios ainda não disponíveis}

Na data desta versão, permanecem pendentes o ciclo integrado com o Odoo para os pares do piloto, a revisão nominal de C1 e do protocolo, a decisão sobre o enquadramento, o dataset final revisado e o congelamento. Não há coleta confirmatória C0--C1 contabilizada. A Tabela~\ref{tab:pending-results} registra o estado das respostas às perguntas de pesquisa; “não executado” não representa zero ocorrência.

\begin{table}[H]
\centering
\caption{Estado dos resultados científicos em 26 de outubro de 2026.}
\label{tab:pending-results}
\small
\begin{tabularx}{\textwidth}{C{0.80cm} L{3.25cm} Y}
\toprule
\textbf{RQ} & \textbf{Estimadores previstos} & \textbf{Estado} \\
\midrule
\hyperlink{rq1}{RQ1} & ASR-proposal, ASR-policy, ASR-state e UACR por C0/C1 & ASR-proposal e ASR-policy medidas no piloto, com 0/36 em cada condição e rodada; ASR-state e UACR não executados; nada confirmatório \\
\hyperlink{rq2}{RQ2} & BTSR, falsos bloqueios, STP, escalonamento, latência e custo & Exatidão das propostas benignas, latência e custo do estágio do modelo medidos no piloto; BTSR no estado, falsos bloqueios e escalonamentos não executados \\
\hyperlink{rq3}{RQ3} & Efeitos por A1, A2 e A3a & Nenhuma família, estratégia, posição ou template produziu proposta comprometida no piloto; teste final não executado \\
\bottomrule
\end{tabularx}
\end{table}

Não é possível afirmar que C1 reduz prompt injection, mantém BTSR no estado do ERP ou possui custo aceitável. O custo será caracterizado, não submetido a hipótese de aceitabilidade. Qwen3, ablação e ataques adaptados à defesa foram retirados da entrega e passam a trabalho futuro; \hyperlink{rq4}{RQ4} e \hyperlink{h5}{H5} não têm resultados a produzir neste estudo. O resumo e a conclusão preservam essa ausência para impedir que números exemplificativos sejam confundidos com observações.

\subsection{Discussão e próximos gates}

A literatura sugere que separar instruções e dados, alinhar ferramentas à tarefa e impor regras em tempo de execução pode reduzir ataques \cite{chen2025struq,jia2025taskshield,li2025drift}. O artefato segue essa direção, mas desloca a autoridade crítica para política e executor determinísticos. Essa escolha procura reduzir a dependência de reconhecer todas as formulações de ataque: mesmo que o LLM proponha outra conta bancária, o estado confiável deve impedir a autorização. O experimento dirá se essa composição reduz violações sem bloquear o processo em excesso.

O piloto altera a leitura do desenho. Com a orientação genérica comum às duas condições, o modelo copiou os valores impressos em todas as execuções adversariais, mesmo diante de autoridade alegada, correção integrada ao documento, ordem direta e texto espaçado. Se a ASR de C0 for nula, a diferença $ASR_{C1}-ASR_{C0}$ não pode ficar abaixo de zero e \hyperlink{h1}{H1} não pode ser satisfeita; como a ordem confirmatória é fixa, \hyperlink{h3}{H3} e \hyperlink{h2}{H2} ficariam sem interpretação confirmatória por um efeito de ataque inexistente. Por isso, há uma proposta, pendente de ratificação do grupo antes do congelamento, de tratar \hyperlink{h1}{H1} como estimativa, com redução absoluta de risco e intervalo, manter \hyperlink{h3}{H3} confirmatória e independente de \hyperlink{h1}{H1}, com a margem inalterada, e decidir \hyperlink{h2}{H2} pela UACR de C0 medida no Odoo. Esta versão registra a proposta, sem adotá-la. O resultado também não demonstra robustez do modelo: os ataques são estáticos, foram calibrados depois da primeira rodada e um único modelo foi avaliado.

Para concluir o piloto, ainda são necessários: ligar o ciclo integrado aos 12 pares com as propostas já geradas, incluindo despacho C0/C1, autorização, supervisão, captura, registro e reset; medir ASR-state, UACR e BTSR; e concluir a revisão independente dos documentos e componentes. O add-on e as provas privadas existentes serão reutilizados, sem expansão de componentes fora desse caminho. O smoke real do modelo e os testes simulados não substituem a competência do agente no piloto. C0 deve cumprir o gate de competência e C1 deve concluir ações legítimas, incluindo contabilização autorizada; a avaliação não pode premiar a simples desativação de toda ação crítica.

Depois do piloto, o grupo deverá documentar discordâncias pareadas, poder, precisão, latência, custo, vazão e completude de logs, além de validar os intervalos nos casos de fronteira. Permanecem 90 pares e três repetições; qualquer revisão da amostra exige decisão explícita antes do teste final, sem corte automático por prazo. O artigo final substituirá a Tabela~\ref{tab:pending-results} por resultados com numeradores, denominadores, missing, diferenças, riscos relativos quando definidos, intervalos e sensibilidades pré-especificadas. Não serão acrescentados blocos de ablação ou replicação local nesta entrega.

O principal risco editorial desta etapa é a falsa sensação de completude produzida por um PDF formatado. O formato institucional já pode ser validado, mas o conteúdo confirmatório não. Por isso, este relatório deve orientar a implementação e a revisão com o orientador, não ser tratado como a submissão final de novembro.
